% New section: bounded low-node deletions for the chi frame.
% This file deliberately leaves the preceding pair argument unchanged.
\section{Uniform nonvanishing after bounded low-node deletion}\label{sec:minor-nv}

For the $\chi$ frame, allowing a bounded number of common nodes to be omitted
removes the restriction $r+2m=3t$ in the preceding pair construction. The
borders remain fixed. We use the rational kernel $M$ of
Section~\ref{sec:construction}. For $J\subseteq[0,r)$ put
\[
 X'=\{r+2i:0\le i<m\},\qquad Y'=\{r+2i+1:0\le i<m\},
\]
and put $\xi_\chi=\tfrac14L(2,\chi_{-3})$. Let $\Phi_J(x)$ have rows $J\sqcup X'$, columns $J\sqcup Y'$, and entries
$M[a,b]+x\mathbf1[a=b\in J]$. Write $P_J(x)=\det\Phi_J(x)$.

\begin{theorem}\label{thm:minor-nv}
Let $\ell=6t+1$ be prime, and suppose
\begin{equation}\label{eq:minor-room}
 m\ge1,\qquad K:=2t+1\le r,\qquad r+2m\le3t.
\end{equation}
If $x\in\Q$ and $\ell^2\nmid\den x$, there is a set $I\subseteq[0,K)$,
$|I|\le3$, such that, for $J=[0,r)\setminus I$,
\[
 P_J(x)\ne0.
\]
\end{theorem}

The proof bounds the dimension of a kernel, rather than asserting simplicity
of its eigenvalues. Conditions~\eqref{eq:minor-room} imply $2m+1\le t$, so
$t\ge3$ and $\ell\ge19$. All occurrences of $2,3,6,18,36$ below are therefore
units modulo $\ell$.

\subsection{Reflected coordinates and the signed moment pencil}
Work first over $k=\F_\ell$. The physical node $b$ and its reflected coordinate
$y=\ell-b$ will be kept distinct. Put
\[
 S=3t+1,\qquad g=\ell-r-2m,\qquad d=K-m,\qquad h=r-K.
\]
Then $1<S\le g$, $3\le d<K$, $d\le g$, and
$K+g+2m+h=\ell$. For a reflected coordinate $y$ define
\[
 c_y=\bigl((a+y)^{-1}\bigr)_{0\le a<K},\qquad
 N=\sum_{y=1}^{g}\epsilon_y c_yc_y^{\top},\qquad
 \epsilon_y=\begin{cases}1&y<S,\\-2&y\ge S.\end{cases}
\]
Reversing the order of each physical border gives the reflected border
coordinates $g+2,g+4,\ldots,g+2m$ and
$g+1,g+3,\ldots,g+2m-1$. Denote their column matrices by $W_X,W_Y$, respectively,
and set
\[
 \Gamma=W_X^{\top}W_Y,\qquad
 B_z=\begin{pmatrix}N-zI_K&W_Y\\W_X^{\top}&0\end{pmatrix}.
\]
Proposition~\ref{prop:delta} gives $\det\Gamma\ne0$: its range condition is
$g+K+2m=\ell-r+K\le\ell$.

Introduce the polynomials
\begin{align*}
 D(y)&=\prod_{a=0}^{K-1}(y+a),&
 B_X(y)&=\prod_{i=0}^{m-1}(y-g-2i-2),\\
 B_Y(y)&=\prod_{i=0}^{m-1}(y-g-2i-1),& R(y)&=B_X(y)B_Y(y).
\end{align*}
For $0\le i<d$ define residue-basis matrices by
\begin{equation}\label{eq:minor-residues}
 (V_X)_{a i}=\frac{B_X(-a)(-a)^i}{D'(-a)},\qquad
 (V_Y)_{a i}=\frac{B_Y(-a)(-a)^i}{D'(-a)}.
\end{equation}
All displayed denominators and all $B_X(-a),B_Y(-a)$ are nonzero.
Partial fractions give, for $\deg p<d$,
\begin{equation}\label{eq:minor-pf}
 \sum_{a=0}^{K-1}\frac{B_X(-a)p(-a)}{D'(-a)(y+a)}
       =\frac{B_X(y)p(y)}{D(y)},
\end{equation}
and likewise for $Y$. Thus the columns of $V_X$ lie in
$\ker W_X^{\top}$. They are independent, since evaluation of a polynomial
of degree $<d\le K$ at the $K$ distinct points $-a$ is injective.
Since $\Gamma$ is invertible, $W_X,W_Y$ have rank $m$; consequently $V_X,V_Y$
are bases of $\ker W_X^{\top},\ker W_Y^{\top}$.

The compressed matrix is the symmetric Hankel pencil
\[
 A_z=V_Y^{\top}(N-zI_K)V_X=G-zH,
\]
whose entries, indexed by $0\le i,j<d$, are
\begin{align}
 H_{ij}&=\sum_{u\in\{0,-1,\ldots,1-K\}}
             \frac{R(u)}{D'(u)^2}u^{i+j},\label{eq:minor-H}\\
 G_{ij}&=\sum_{y=1}^{g}\epsilon_y\frac{R(y)}{D(y)^2}y^{i+j}.
 \label{eq:minor-G}
\end{align}
These identities follow from~\eqref{eq:minor-residues} and
\eqref{eq:minor-pf}. In particular the same polynomial coordinates describe
the left and right kernels; this fact will be needed for principal deletion.

\subsection{A uniform corank bound}
\begin{lemma}\label{lem:minor-corank}
For every $z\in k$, $\dim_k\ker A_z\le3$.
\end{lemma}
\begin{proof}
Identify a coefficient vector in $\ker A_z$ with a polynomial $p$, $\deg p<d$.
Equations~\eqref{eq:minor-H}--\eqref{eq:minor-G} say that its first $d$
weighted moments vanish. To use these conditions, complete the root sets by
\[
 E(y)=\prod_{j=1}^{g}(y-j),\qquad
 C(y)=\prod_{j=g+2m+1}^{\ell-K}(y-j).
\]
The last product is $1$ when $h=0$. The roots of $D,E,R,C$ partition $k$, so
\begin{equation}\label{eq:minor-fullfield}
 DERC=y^\ell-y,\qquad (y^\ell-y)'=-1.
\end{equation}
Let $T=DE$. On its simple roots define
\[
 w_z(v)=\begin{cases}
 -zR(v)/D'(v)^2&D(v)=0,\\
 \epsilon_v R(v)/D(v)^2&E(v)=0,
 \end{cases}
\]
and form the polynomial, linear in $p$,
\[
 B_p(y)=\sum_{T(v)=0}w_z(v)p(v)\frac{T(y)}{y-v},\qquad
 F_p(y)=D(y)^2C(y)B_p(y).
\]
The first $d$ moments are zero. Expanding
$B_p/T=\sum_v w_z(v)p(v)/(y-v)$ formally at infinity therefore gives
\[
 \deg B_p\le K+g-d-1=g+m-1,\qquad \deg F_p\le\ell+d-1.
\]
Equivalently these degree bounds follow by comparing the first $d$ leading
coefficients in the finite polynomial sum defining $B_p$.
Evaluating at a root of $T$, and differentiating
\eqref{eq:minor-fullfield}, gives the exact identities
\begin{equation}\label{eq:minor-interpolation}
 \begin{array}{ll}
 F_p(u)=F_p'(u)=0,\quad F_p''(u)=2zp(u)&(D(u)=0),\\
 F_p(j)=-\epsilon_jp(j)&(1\le j\le g),\\
 F_p(j)=0&(C(j)=0).
 \end{array}
\end{equation}
For example $B_p(u)=-zR(u)E(u)p(u)/D'(u)$ at a root $u$ of $D$;
then $F_p''(u)=2D'(u)^2C(u)B_p(u)=2zp(u)$ by
$D'(u)E(u)R(u)C(u)=-1$. The high-node identity follows in the same way.

Consider the finite-difference Wronskian
\[
 W_p(y)=F_p(y+1)p(y)-F_p(y)p(y+1).
\]
At each low adjacency $-K+1,\ldots,-1$, its value and first derivative
vanish. Its second derivative also vanishes, by the equality
$2zp(y+1)p(y)-2zp(y)p(y+1)=0$. Hence these are triple roots.
At the high adjacencies $1,\ldots,S-2$ and $S,\ldots,g-1$ its value
vanishes because the two weights agree. It vanishes at every root of $C$:
for the last such root, $\ell-K$, the next point is the first low root
$\ell-K+1$. Thus
\begin{equation}\label{eq:minor-fixedfactor}
 Q_0(y):=\prod_{a=1}^{K-1}(y+a)^3
       \prod_{j=1}^{S-2}(y-j)\prod_{j=S}^{g-1}(y-j)C(y)
       \quad\hbox{divides } W_p(y).
\end{equation}
The supports are disjoint, and
\[
 \begin{gathered}
 \deg Q_0=3(K-1)+(g-2)+h=\ell+2d-5,
 \\ Q_0(S-1)\ne0,\qquad W_p(S-1)=3p(S)p(S-1).
 \end{gathered}
\]

We also need that $W_p\ne0$ for nonzero $p$ in the kernel.
If $W_p=0$, reduce $F_p/p=A/B$ to coprime polynomials with $B$ monic.
Then $A(y+1)B(y)=A(y)B(y+1)$ implies $B\mid B(y+1)$, and hence
$B(y+1)=B(y)$. But $\deg B\le\deg p<\ell$; a nonconstant polynomial
of degree $e<\ell$ has a nonzero leading term of degree $e-1$ in its
first difference. Thus $B$ is constant. The polynomial quotient $F_p/p$
is translation invariant and has degree $<2\ell$. Its positive degree,
if any, must be $\ell$ by the same leading-coefficient argument.
Subtracting a multiple of $y^\ell-y$ therefore shows
\[
 F_p=p\bigl(a(y^\ell-y)+b\bigr).
\]
Some low node $u$ has $p(u)\ne0$, since $\deg p<K$. The first two
identities in~\eqref{eq:minor-interpolation} force $b=0$ and $a=0$.
Then the high-node identity forces $p$ to vanish at $g>\deg p$ points,
a contradiction.

Finally, if the kernel had dimension at least four, its three linear
conditions
\[
 [y^{d-1}]p=[y^{d-2}]p=0,\qquad p(S)=0
\]
would admit a nonzero solution. This solution has $\deg p\le d-3$.
The leading terms in the two products defining $W_p$ cancel, so
$\deg W_p\le\deg F_p+\deg p-1\le\ell+2d-5=\deg Q_0$.
By~\eqref{eq:minor-fixedfactor}, $W_p$ is a scalar multiple of $Q_0$;
its value at $S-1$ makes that scalar zero. This contradicts the preceding
paragraph and proves the lemma.
\end{proof}

\subsection{Selecting a nonzero principal minor}
\begin{lemma}\label{lem:minor-rescue}
For every $z\in k$, deleting at most three low principal coordinates from
$B_z$ leaves a nonsingular matrix.
\end{lemma}
\begin{proof}
Let $\mathcal P_z$ be the polynomial kernel of $A_z$, and put
$k_0=\dim\mathcal P_z\le3$. A right kernel vector of $B_z$ has low
part $v\in\ker W_X^{\top}$, hence $v=V_Xp$ uniquely. Its remaining
equation is $(N-zI)v\in\operatorname{im}W_Y$, which is equivalent to
$V_Y^{\top}(N-zI)V_Xp=0$. The border coordinates are then unique,
since $\Gamma$ is invertible. Conversely every $p\in\mathcal P_z$
lifts to such a kernel vector. The transpose has the corresponding
parametrization with low part $V_Yp$; the polynomial space is the same
because $A_z$ is symmetric.

Evaluation $p\mapsto(p(-a))_{a<K}$ is injective on $\mathcal P_z$.
Choose $k_0$ coordinates on which it is an isomorphism, and let $I$ be
the corresponding low indices. By~\eqref{eq:minor-residues}, restriction
of either full kernel to $I$ is an isomorphism: it differs from polynomial
evaluation only by nonzero diagonal factors.

Here is the elementary principal-minor criterion being used. Suppose a
square matrix has injective restriction of its right kernel to $I$ and
surjective restriction of its left kernel to $I$. A kernel vector of
the complementary principal matrix, extended by zero on $I$, is sent by
the full matrix to a vector supported on $I$. Pairing with left kernel
vectors shows that this image is zero. Right-kernel injectivity then
shows the extended vector is zero. Thus the complementary principal
matrix is nonsingular. Apply this criterion to $B_z$.
\end{proof}

\subsection{Return to the physical determinant}
Suppose first $\ell\nmid\den x$, and write $\bar x$ for its reduction.
Let
\[
 c_0=\sum_{y=1}^{3t}y^{-2}-2\sum_{y=3t+1}^{4t}y^{-2},\quad
 \kappa=-\frac1{18},\quad \tau_0=\frac1{36},\quad h_0=\frac{c_0}{2}.
\]
The three-level entry calculation of Section~\ref{ssec:nvother} applies
under~\eqref{eq:minor-room}; it does not require equality $r+2m=3t$.
The physical absent middle nodes are $[r+2m,S)$, which reflect to $[S,g]$;
the top nodes reflect to $[1,S)$. Their weighted Gram matrix is therefore
\[
 N_{\rm phys}=\sum_{y=S}^{g}c_yc_y^{\top}
              -\frac12\sum_{y=1}^{S-1}c_yc_y^{\top}=-\frac12 N.
\]
Consequently the physical reduced low block is exactly
\begin{equation}\label{eq:minor-physical}
 (h_0+\bar x/\kappa)I+N_{\rm phys}
       =-\frac12\bigl(N-zI\bigr),\qquad z=c_0-36\bar x.
\end{equation}
Both border orders are reversed simultaneously to obtain $W_X,W_Y$.
Scaling a low block by a nonzero scalar $c$ preserves bordered
nonsingularity, since
\[
 \begin{pmatrix}cI&0\\0&I\end{pmatrix}
 \begin{pmatrix}M_0&W_Y\\W_X^{\top}&0\end{pmatrix}
 \begin{pmatrix}I&0\\0&c^{-1}I\end{pmatrix}
 =\begin{pmatrix}cM_0&W_Y\\W_X^{\top}&0\end{pmatrix}.
\]
This also holds after any low principal restriction. Thus
Lemma~\ref{lem:minor-rescue} supplies a nonzero physical bordered minor.

For completeness, this minor transfers to the rational matrix without
assuming that the retained common set contains every low node.
Start with the original geometry and multiply its low rows and columns
by $1$, its middle rows and columns by $\ell$, and negate the matrix.
For any retained low set $U\subseteq[0,K)$ its reduced principal
restriction has the same Schur identity as the full matrix. Writing
$M_{\rm phys}$ for the left-hand side of~\eqref{eq:minor-physical},
\begin{equation}\label{eq:minor-schur-factor}
 \det Z_U=(1/2)^{r-K+m}(1/18)^{|U|}
 \det\begin{pmatrix}
 M_{\rm phys}[U,U]&W_Y[U,*]\\
 W_X[U,*]^{\top}&0
 \end{pmatrix}.
\end{equation}
Indeed the invertible middle common block is unchanged by low restriction;
the harmonic identity restricts entrywise, and the border row and column
operations are unchanged. Every prefactor is a unit. A zero rational
principal determinant would remain zero after the rational diagonal
scalings and after reduction, contradicting~\eqref{eq:minor-schur-factor}.
Taking $U=[0,K)\setminus I$ proves the theorem in this case.

It remains to treat $\ell\mid\den x$ but $\ell^2\nmid\den x$.
Coprimality of the reduced numerator and denominator shows that
$u=\overline{\ell x}$ is nonzero. Use the undeleted common set $J=[0,r)$.
Multiply \emph{every} row of $\Phi_J(x)$ by $\ell$ and additionally
multiply each middle common column by $\ell$. With rows and columns
ordered as low, middle common, border, the resulting integral matrix
reduces to
\[
 \begin{pmatrix}
 uI_K&0&-\tfrac16W_Y\\
 -\tfrac16W_C^{\top}&-\tfrac12I_{r-K}&0\\
 -\tfrac16W_X^{\top}&0&0
 \end{pmatrix}.
\]
Here $W_C$ consists of the middle common columns. In verifying the zero
blocks one uses integrality of distinct middle--middle entries, rather
than merely integrality after multiplication by $\ell^2$; it follows
directly from the off-diagonal divided-difference formula, whose interval
contains neither singular index $K$ nor $S$.
Reorder the middle common block last. Its determinant is $(-1/2)^{r-K}$,
and the remaining bordered determinant is
\[
 \det\begin{pmatrix}uI_K&-\tfrac16W_Y\\-\tfrac16W_X^{\top}&0\end{pmatrix}
       =u^{K-m}(-1/36)^m\det\Gamma\ne0.
\]
Thus $P_{[0,r)}(x)\ne0$. This proves
Theorem~\ref{thm:minor-nv} in the remaining case.

\subsection{Content, decay, and the original degree budget}
The enlarged family consists of all $P_J$ with $|[0,r)\setminus J|\le3$.
For admissible $r\le8m\le r+9$ it retains the constants
\[
 \Lambda=\frac{10339}{3072},\qquad D=3.7251,\qquad c=6.6510.
\]
To see the content assertion, use the entrywise tropical denominator
bound from Section~\ref{sec:content}, before any determinant cancellation.
Any assignment for a retained principal submatrix extends to one for
the original matrix by fixing the omitted common coordinates. All
weights are nonnegative, so the smaller assignment value cannot be
larger. Write $d_{r,m}$ for this common integer denominator. It clears every $P_J$, and
its logarithm is at most $(\Lambda+o(1))r^2$, uniformly in $J$.

The energy estimates of Section~\ref{sec:decay} already apply to every
subset of common nodes, with the borders retained. If $k\le3$ common
nodes are forced to be absent, expansion at $\xi_\chi+v$ has terms
\[
 v^{r-k-|A|}\det\bigl(\omega\Lw_{F_\chi}[A\sqcup X',A\sqcup Y']\bigr),
 \qquad A\subseteq J.
\]
For $|v|\le e^{-cr}$, use
$|v|^{r-k-|A|}\le e^{-cr(r-k-|A|)}$. Replacing the original exponent
$r-|A|$ by $r-k-|A|$ therefore costs at most $e^{crk}$ in the termwise
upper bound, also at $v=0$. This is
a logarithmic loss of at most $3cr=O(r)$, and the number of terms is
still at most $2^r$. It is absorbed by the existing $o(r^2)$ errors.
Hence $|P_J(\xi_\chi+v)|\le\exp(-(D-o(1))r^2)$ on the same window,
uniformly in $J$; no comparison of possibly cancelling determinants is
needed.

Finally $\deg P_J\le|J|\le r$. We keep the declared degree budget equal
to the original $r$ for every member, so $q^r$ times the common clearing
denominator times $P_J(p/q)$ is an integer. Neither the degree scale nor
the exponential window changes when the nonzero member is selected.
